{\color{gray}\hrule}
\begin{center}
\section{Results}
% \textbf{Small description}
\bigskip
\end{center}
{\color{gray}\hrule}
\begin{multicols}{2}

The primary behavioral evaluation focused on the classification strategies of the three networks (standard, biomimetic, anti-biomimetic) when presented with conflicting visual cues. Using the cue-conflict dataset,\cite{geirhos2018a} we quantified the extent to which each regimen developed a "shape bias"—a characteristic typically associated with human visual development but often lacking in standard DNNs.

\subsection{Shape vs. Texture Bias Analysis}
The results of the cue-conflict test reveal significant differences in how the temporal order of sensory input quality influences the final classification strategy of the networks.

\begin{itemize}
    \item \textbf{Standard Regimen:} Consistent with previous literature, the control network trained exclusively on high-resolution images exhibited a strong \textit{texture bias}. When presented with an image containing the texture of one class (e.g., a "clock") mapped onto the shape of another (e.g., a "bottle"), the network predominantly classified the object based on the local texture.
    
    \item \textbf{Biomimetic Regimen:} The network that began training with blurred, achromatic inputs showed a marked increase in \textit{shape bias} compared to the standard regimen. By the end of the 120 epochs, this network demonstrated a higher percentage of shape-consistent decisions, suggesting that the initial period of low-spatial-frequency information forced the network to learn global structural features rather than local high-frequency textures.
    
    \item \textbf{Anti-Biomimetic Regimen:} In contrast, the network subjected to the reversed order (high-resolution followed by blurry images) failed to develop a robust shape bias. Interestingly, the transition to degraded inputs in the second half of training did not retroactively instill a shape-centric strategy; instead, it resulted in a degradation of overall classification accuracy while maintaining a residual reliance on the features learned during the first high-resolution phase.
\end{itemize}

\subsection{Quantification of Bias}
The degree of shape bias $S$ was calculated as the proportion of shape-consistent decisions $D_{shape}$ relative to the total number of correct category decisions (shape-consistent + texture-consistent, $D_{texture}$):

\begin{equation}
S = \frac{D_{shape}}{D_{shape} + D_{texture}}
\end{equation}

[Inference] Preliminary analysis indicates that the Biomimetic regimen achieves a significantly higher $S$ value than both the Standard and Anti-Biomimetic regimens, supporting the hypothesis that the developmental "starting small" approach is crucial for human-like visual processing.

\end{multicols}